%!TEX root = ../main.tex
% GLOSARIO: términos propios del trabajo.

\newglossaryentry{cuidadorinformal}{
    name={cuidador informal},
    description={Persona no profesional que brinda asistencia y acompañamiento
    habitual a una persona con discapacidad, por lo general sin recibir un pago
    ni contar con formación previa especializada}
}

\newglossaryentry{cuidadorcontratado}{
    name={cuidador contratado},
    description={Persona que cuida de forma remunerada, contratada por una
    familia o empleada por una institución}
}

\newglossaryentry{grupocuidado}{
    name={grupo de cuidado},
    description={Conjunto de personas que participan del cuidado dentro de la
    plataforma: el dueño, los cuidadores y los observadores}
}

\newglossaryentry{dueno}{
    name={dueño del grupo},
    description={Usuario que crea la cuenta, carga a las personas asistidas,
    suma integrantes y paga el plan, que cubre a todo su grupo}
}

\newglossaryentry{observador}{
    name={observador},
    description={Integrante de un grupo que puede ver lo que se registra sobre
    las personas que tiene asignadas, sin poder modificarlo}
}

\newglossaryentry{personaasistida}{
    name={persona asistida},
    description={Persona con discapacidad cuyo cuidado se organiza en la
    plataforma}
}

\newglossaryentry{recomendadorhibrido}{
    name={sistema de recomendación híbrido},
    description={Sistema que combina el filtrado basado en contenido con el
    filtrado colaborativo para ordenar sus sugerencias}
}

\newglossaryentry{filtradocontenido}{
    name={filtrado basado en contenido},
    description={Técnica de recomendación que compara las características del
    usuario con las de cada ítem, sin depender del comportamiento de otros
    usuarios}
}

\newglossaryentry{filtradocolaborativo}{
    name={filtrado colaborativo},
    description={Técnica de recomendación que sugiere lo que resultó útil a
    usuarios con perfiles parecidos}
}

\newglossaryentry{arranquefrio}{
    name={arranque en frío},
    description={Situación en la que un sistema de recomendación no tiene datos
    suficientes sobre un usuario o un ítem nuevo para calcular similitudes}
}

\newglossaryentry{valoracionexplicita}{
    name={valoración explícita},
    description={Indicación de un cuidador sobre si un recurso recomendado le
    sirvió o no}
}

\newglossaryentry{sobrecarga}{
    name={sobrecarga del cuidador},
    description={Conjunto de efectos físicos, emocionales y sociales que surgen
    de la acumulación de estresores del cuidado cuando los recursos de
    afrontamiento no alcanzan}
}

\newglossaryentry{prototipo}{
    name={prototipo},
    description={Versión preliminar de un sistema construida para validar una
    idea, un flujo o una interfaz}
}

\newglossaryentry{migracion}{
    name={migración},
    plural={migraciones},
    description={Archivo versionado que describe un cambio en el esquema de la
    base de datos y permite aplicarlo o revertirlo}
}

\newglossaryentry{exposicion}{
    name={exposición al riesgo},
    description={Producto de la probabilidad de un riesgo por su impacto, que
    se usa para ordenar la matriz de riesgos}
}
